\documentclass[aspectratio=169,10pt]{beamer}

\usepackage[utf8]{inputenc}
\usepackage[T1]{fontenc}
\usepackage[spanish,es-nodecimaldot]{babel}
\usepackage{helvet}
\usepackage{microtype}
\usepackage{booktabs}
\usepackage{tabularx}
\usepackage{tikz}
\usetikzlibrary{arrows.meta,positioning,calc,shapes.geometric,fit}
\renewcommand{\familydefault}{\sfdefault}

\definecolor{Navy}{HTML}{0A2342}
\definecolor{Blue}{HTML}{176B87}
\definecolor{Cyan}{HTML}{3BB7C4}
\definecolor{Mint}{HTML}{DDF5F1}
\definecolor{Gold}{HTML}{F4B942}
\definecolor{Coral}{HTML}{F26B5B}
\definecolor{Ink}{HTML}{18212F}
\definecolor{Slate}{HTML}{5D6A7A}
\definecolor{Fog}{HTML}{F2F5F7}
\definecolor{White}{HTML}{FFFFFF}

\setbeamercolor{normal text}{fg=Ink,bg=White}
\setbeamercolor{structure}{fg=Blue}
\setbeamercolor{frametitle}{fg=Navy,bg=White}
\setbeamerfont{frametitle}{series=\bfseries,size=\Large}
\setbeamerfont{title}{series=\bfseries,size=\fontsize{28}{32}}
\setbeamerfont{subtitle}{size=\fontsize{13}{16}}
\setbeamertemplate{navigation symbols}{}
\setbeamertemplate{itemize item}{\color{Cyan}\raise0.15ex\hbox{\tiny$\blacktriangleright$}}
\setbeamertemplate{itemize subitem}{\color{Cyan}\raise0.1ex\hbox{\tiny$\bullet$}}
\setlength{\leftmargini}{1.25em}
\setlength{\leftmarginii}{1.1em}

\setbeamertemplate{frametitle}{%
  \vspace{0.35cm}\insertframetitle\par
  \vspace{0.08cm}{\color{Cyan}\rule{1.35cm}{1.8pt}}
}
\setbeamertemplate{footline}{%
  \leavevmode
  \hbox{%
  \begin{beamercolorbox}[wd=.83\paperwidth,ht=2.7ex,dp=1.2ex,leftskip=.55cm]{normal text}%
    \color{Slate}\scriptsize KALEIDO FLOWTWIN \quad | \quad PROPUESTA PARA CONVERSACION
  \end{beamercolorbox}%
  \begin{beamercolorbox}[wd=.17\paperwidth,ht=2.7ex,dp=1.2ex,rightskip=.55cm plus1fil]{normal text}%
    \color{Slate}\scriptsize \insertframenumber/\inserttotalframenumber
  \end{beamercolorbox}}
}

\newcommand{\eyebrow}[1]{\textcolor{Cyan}{\bfseries\scriptsize\MakeUppercase{#1}}}
\newcommand{\metric}[2]{%
  \begin{minipage}[t]{\linewidth}
  {\color{Navy}\bfseries\Large #1}\par\vspace{1mm}
  {\color{Slate}\small #2}
  \end{minipage}}
\newcommand{\tagbox}[2][Mint]{%
  \begingroup\setlength{\fboxsep}{4pt}\colorbox{#1}{\scriptsize\bfseries\color{Ink}\strut #2}\endgroup}

\title{Kaleido FlowTwin}
\subtitle{De trazabilidad a anticipacion operativa, sin depender de un historico rico}
\author{Álvaro Schwiedop Souto -- Industrial AI \& Predictive Quality Lead, EVOCON Solutions GmbH}
\date{Conversacion tecnica -- 15 julio 2026}

\begin{document}

{\setbeamertemplate{footline}{}
\begin{frame}[plain]
  \begin{tikzpicture}[remember picture,overlay]
    \fill[Navy] (current page.south west) rectangle (current page.north east);
    \fill[Cyan] (current page.south west) rectangle ([xshift=0.18cm]current page.north west);
    \fill[Blue,opacity=.42] ([xshift=9.7cm,yshift=-7.5cm]current page.north west) circle (4.1cm);
    \fill[Cyan,opacity=.18] ([xshift=13.6cm,yshift=-1.1cm]current page.north west) circle (3.5cm);
  \end{tikzpicture}
  \vspace{0.7cm}
  \eyebrow{Propuesta evidence-first}\par\vspace{0.45cm}
  {\color{White}\usebeamerfont{title}\inserttitle}\par\vspace{0.22cm}
  {\color{White}\usebeamerfont{subtitle}\insertsubtitle}\par
  \vspace{1.15cm}
  \begin{minipage}{0.68\linewidth}
  {\color{Mint}\large Riesgo de desviacion, tiempo restante y escenarios
  sobre los eventos que Kaleido ya captura.}
  \end{minipage}
  \vfill
  {\color{White}\small \insertauthor}\hfill
  {\color{Mint}\small \insertdate}
\end{frame}}

\begin{frame}{La restriccion revela la oportunidad}
  \eyebrow{Lo que hemos oido de Kaleido}
  \vspace{0.55cm}
  \begin{tikzpicture}
    \node[fill=Fog,rounded corners=7pt,inner sep=15pt,text width=.87\linewidth,anchor=west] (q)
      {\color{Navy}\Large\bfseries ``No contamos con un historico de datos de operativas muy rico.''};
    \draw[Cyan,line width=3pt] (q.north west) -- (q.south west);
  \end{tikzpicture}
  \vspace{0.65cm}

  \begin{columns}[T,onlytextwidth]
    \column{.47\textwidth}
      {\bfseries\color{Navy}La oportunidad}\par\vspace{2mm}
      \small Crear la memoria operativa que multiplica el valor de Trace Port y mejora con cada nueva operacion.
    \column{.47\textwidth}
      {\bfseries\color{Navy}La respuesta}\par\vspace{2mm}
      \small Arrancar con process mining, baselines y simulacion; incorporar
      Event-JEPA a medida que crece el activo de datos.
  \end{columns}
\end{frame}

\begin{frame}{De Trace Port a inteligencia operativa predictiva}
\eyebrow{Producto y vision}
  \vspace{0.35cm}
  \begin{columns}[T,onlytextwidth]
    \column{.47\textwidth}
      \begin{tikzpicture}
        \node[fill=Mint,draw=Cyan!60,rounded corners=6pt,inner sep=11pt,text width=.88\linewidth] {
          {\bfseries\color{Blue}MVP en 4--6 semanas}\par\vspace{2mm}
          \small Estado operativo, tiempo restante, riesgo de desviacion,
          cuellos de botella y escenarios sobre una operacion real.};
      \end{tikzpicture}
    \column{.47\textwidth}
      \begin{tikzpicture}
        \node[fill=Blue!12,draw=Blue!55,rounded corners=6pt,inner sep=11pt,text width=.88\linewidth] {
          {\bfseries\color{Navy}Vision de plataforma}\par\vspace{2mm}
          \small Un sistema que aprende entre terminales, simula planes y
          convierte Trace Port en un producto prescriptivo de alto margen.};
      \end{tikzpicture}
  \end{columns}
  \vspace{0.7cm}
  \centering
  \tagbox{Kaleido ya posee workflow, usuarios y data exhaust: la base de una ventaja acumulativa}
\end{frame}

\begin{frame}{Kaleido ya tiene el activo que hace viable la idea}
  \eyebrow{Encaje, no sustitucion}
  \vspace{0.35cm}
  \begin{columns}[T,onlytextwidth]
    \column{.31\textwidth}
      \begin{tikzpicture}
      \node[fill=Fog,rounded corners=6pt,inner sep=6pt,text width=.84\linewidth,minimum height=2.72cm] {
        {\bfseries\color{Navy}Trace Port}\par\smallskip
        \small Proyectos, packing lists, turnos, eventos, equipos, fotos,
        incidencias, historico e integracion.\par\medskip
        \tagbox{Primer modulo}};
      \end{tikzpicture}
    \column{.31\textwidth}
      \begin{tikzpicture}
      \node[fill=Fog,rounded corners=6pt,inner sep=6pt,text width=.84\linewidth,minimum height=2.72cm] {
        {\bfseries\color{Navy}Shipping Board}\par\smallskip
        \small Visibilidad compartida de expediciones y recepciones, documentos,
        alertas y estadisticas.\par\medskip
        \tagbox[Gold!25]{Extension}};
      \end{tikzpicture}
    \column{.31\textwidth}
      \begin{tikzpicture}
      \node[fill=Fog,rounded corners=6pt,inner sep=6pt,text width=.84\linewidth,minimum height=2.72cm] {
        {\bfseries\color{Navy}Freight Intelligence}\par\smallskip
        \small Seguimiento de contenedores, eventos de carriers, alertas e
        integracion.\par\medskip
        \tagbox[Gold!25]{Extension}};
      \end{tikzpicture}
  \end{columns}
  \vspace{0.1cm}
  \centering\small\color{Slate}
  FlowTwin convierte trazabilidad en anticipacion sin rehacer el producto existente.
\end{frame}

\begin{frame}{Primer caso: anticipar una operacion de Trace Port}
  \eyebrow{Pregunta de negocio}
  \vspace{0.35cm}
  \begin{tikzpicture}
    \node[fill=Navy,rounded corners=7pt,inner sep=13pt,text width=.9\linewidth,text=White] {
      \large\bfseries Dado el estado actual, el plan, los recursos y el contexto:
      ¿terminaremos a tiempo y que puede hacerse antes de que sea tarde?};
  \end{tikzpicture}
  \vspace{0.55cm}

  \begin{columns}[T,onlytextwidth]
    \column{.25\textwidth}
      {\bfseries\color{Blue}Tiempo restante}\par\smallskip
      \small Estimacion e intervalo hasta fin o milestone.
    \column{.25\textwidth}
      {\bfseries\color{Blue}Riesgo 2/4/8 h}\par\smallskip
      \small Probabilidad calibrada de exceder el plan.
    \column{.25\textwidth}
      {\bfseries\color{Blue}Cuello de botella}\par\smallskip
      \small Donde se acumula espera y que evidencia lo explica.
    \column{.25\textwidth}
      {\bfseries\color{Blue}Escenarios}\par\smallskip
      \small Comparar solo acciones aprobadas; humano decide.
  \end{columns}
  \vspace{0.6cm}
  \small\color{Slate}Primera unidad de valor: una alerta suficientemente temprana para una decision real.
\end{frame}

\begin{frame}{Escalera de valor: del dato a la optimizacion}
\eyebrow{Valor desde el primer ciclo}
  \vspace{0.55cm}
  \begin{tikzpicture}[node distance=.12cm]
    \tikzset{cstep/.style={rounded corners=5pt,minimum width=2.58cm,minimum height=1.25cm,text width=2.20cm,align=center,font=\small\bfseries}}
    \node[cstep,fill=Fog,draw=Slate!35] (s1) {1. Espejo operativo};
    \node[cstep,fill=Mint,draw=Cyan!55,right=of s1] (s2) {2. Insight y bottlenecks};
    \node[cstep,fill=Cyan!22,draw=Cyan!70,right=of s2] (s3) {3. Prediccion 2/4/8 h};
    \node[cstep,fill=Gold!24,draw=Gold!80,right=of s3] (s4) {4. Escenarios what-if};
    \node[cstep,fill=Blue!15,draw=Blue!65,right=of s4] (s5) {5. Optimizacion continua};
    \foreach \a/\b in {s1/s2,s2/s3,s3/s4,s4/s5}
      \draw[-{Latex[length=2mm]},line width=1.1pt,color=Slate!70] (\a) -- (\b);
  \end{tikzpicture}
  \vspace{0.7cm}
  \begin{columns}[T,onlytextwidth]
    \column{.48\textwidth}
      \tagbox{Valor inmediato}\par\smallskip
      \small Process mining, tiempos y variantes convierten eventos dispersos en una vision operativa comun.
    \column{.48\textwidth}
      \tagbox[Gold!25]{Activo acumulativo}\par\smallskip
      \small Cada operacion alimenta mejores predicciones, simulaciones y productos para nuevos clientes.
  \end{columns}
\end{frame}

\begin{frame}{Arquitectura: del evento a una recomendacion auditable}
  \eyebrow{Read-only por diseno}
  \vspace{0.25cm}
  \centering
  \begin{tikzpicture}[node distance=.43cm and .48cm,scale=.69,transform shape]
    \tikzset{box/.style={rounded corners=4pt,minimum height=.82cm,align=center,font=\scriptsize\bfseries,inner xsep=7pt},
      arr/.style={-{Latex[length=2.1mm]},line width=1.1pt,color=Slate!75}}
    \node[box,fill=Navy,text=White,text width=2.35cm] (src) {Trace Port\\ERP / TOS / meteo};
    \node[box,fill=Fog,draw=Slate!35,text width=2.55cm,right=of src] (evt) {Eventos versionados\\calidad temporal};
    \node[box,fill=Mint,draw=Cyan!55,text width=2.35cm,right=of evt] (state) {Estado operativo\\a tiempo $t$};
    \node[box,fill=Cyan!18,draw=Cyan!70,text width=2.35cm,right=of state] (model) {Baseline + secuencial\\+ simulador};
    \node[box,fill=Gold!24,draw=Gold!80,text width=2.35cm,right=of model] (risk) {Riesgo + tiempo\\+ intervalo};
    \node[box,fill=Blue,text=White,text width=2.2cm,right=of risk] (ui) {Trace Port / API\\humano decide};
    \draw[arr] (src) -- (evt);
    \draw[arr] (evt) -- (state);
    \draw[arr] (state) -- (model);
    \draw[arr] (model) -- (risk);
    \draw[arr] (risk) -- (ui);
    \node[below=.55cm of state,fill=Coral!12,draw=Coral!55,rounded corners=4pt,inner sep=5pt,font=\scriptsize] (gate) {Fuga, OOD, abstencion, versionado};
    \draw[-{Latex[length=2mm]},dashed,color=Coral!75] (gate) -- (state);
    \draw[-{Latex[length=2mm]},dashed,color=Coral!75] (gate.east) -- (risk.south);
  \end{tikzpicture}
  \vspace{0.7cm}

  \begin{columns}[T,onlytextwidth]
    \column{.32\textwidth}\metric{2 tiempos}{event\_time e ingest\_time}
    \column{.32\textwidth}\metric{0 writes}{sin control automatico}
    \column{.32\textwidth}\metric{1 rastro}{dato, modelo y explicacion}
  \end{columns}
\end{frame}

\begin{frame}{El contrato de datos evita una demo enganosa}
  \eyebrow{Minimo viable de evidencia}
  \vspace{0.25cm}
  \begin{columns}[T,onlytextwidth]
    \column{.47\textwidth}
      {\bfseries\color{Navy}Entidades y proceso}\par\vspace{2mm}
      \small
      \begin{itemize}
        \item proyecto, operacion, turno y unidad de carga;
        \item evento, estado, recurso/rol e incidencia;
        \item plan y milestones con version vigente;
        \item inicio, pausa, reanudacion y fin.
      \end{itemize}
    \column{.47\textwidth}
      {\bfseries\color{Navy}Semantica temporal}\par\vspace{2mm}
      \small
      \begin{itemize}
        \item tiempo del hecho y tiempo de llegada;
        \item correcciones, duplicados y eventos tardios;
        \item conocimiento disponible en cada snapshot;
        \item resultado final separado del prefijo.
      \end{itemize}
  \end{columns}
  \vspace{0.4cm}
  \begin{tikzpicture}
    \node[fill=Fog,rounded corners=5pt,inner sep=9pt,text width=.9\linewidth] {
      \small\bfseries\color{Blue}Muestra inicial recomendada: esquema + perfiles agregados + 3--5 operaciones completas y pseudonimizadas.};
  \end{tikzpicture}
\end{frame}

\begin{frame}{Stack de producto: fiable hoy, diferencial manana}
\eyebrow{Arquitectura multimodelo}
  \vspace{0.35cm}
  \begin{columns}[T,onlytextwidth]
    \column{.31\textwidth}
      \begin{tikzpicture}
        \node[fill=Fog,rounded corners=6pt,inner sep=7pt,text width=.84\linewidth,minimum height=3.15cm] {
        {\bfseries\color{Navy}Baselines}\par\smallskip
        \small Reglas, mediana por familia, Kaplan--Meier, quantile boosting.\par\medskip
        \tagbox{Capa de control}};
      \end{tikzpicture}
    \column{.31\textwidth}
      \begin{tikzpicture}
        \node[fill=Mint,rounded corners=6pt,inner sep=7pt,text width=.84\linewidth,minimum height=3.15cm] {
        {\bfseries\color{Navy}PPM / grafos}\par\smallskip
        \small Transformer de prefijos y grafo object-centric para tiempo y siguiente evento.\par\medskip
        \tagbox[Gold!25]{Motor operativo}};
      \end{tikzpicture}
    \column{.31\textwidth}
      \begin{tikzpicture}
        \node[fill=Blue!12,rounded corners=6pt,inner sep=7pt,text width=.84\linewidth,minimum height=3.15cm] {
        {\bfseries\color{Navy}Event-JEPA}\par\smallskip
        \small Predecir la representacion futura a 2/4/8 h, con regularizacion y probes.\par\medskip
        \tagbox{Ventaja EVOCON}};
      \end{tikzpicture}
  \end{columns}
\end{frame}

\begin{frame}{Por que ahora: convergen proceso, JEPA y simulacion}
\eyebrow{Ventana tecnologica 2026}
  \vspace{0.25cm}
  \begin{columns}[T,onlytextwidth]
    \column{.48\textwidth}
      {\bfseries\color{Blue}El SOTA ya permite}\par\smallskip
      \small
      \begin{itemize}
        \item predecir estado latente, no cada detalle;
        \item horizontes directos y jerarquia temporal;
        \item acciones separadas del contexto;
        \item anticolapso, incertidumbre y probes;
        \item adaptacion entre terminales y operativas.
      \end{itemize}
    \column{.48\textwidth}
      {\bfseries\color{Navy}La ventaja de Kaleido + EVOCON}\par\smallskip
      \small
      \begin{itemize}
        \item Trace Port ya vive dentro de la operacion;
        \item Kaleido conoce procesos, clientes y decisiones;
        \item EVOCON aporta world models y predictive operations;
        \item el producto crea datos mejores mientras se usa;
        \item la misma plataforma escala a Shipping Board y Freight Intelligence.
      \end{itemize}
  \end{columns}
  \vspace{0.25cm}
  \centering\tagbox{La oportunidad: convertir visibilidad en anticipacion y anticipacion en optimizacion}
\end{frame}

\begin{frame}{Prueba de valor: metricas que desbloquean compra}
\eyebrow{Rendimiento tecnico + impacto operativo}
  \vspace{0.2cm}
  \begin{columns}[T,onlytextwidth]
    \column{.47\textwidth}
      {\bfseries\color{Navy}Calidad predictiva}\par\smallskip
      \small
      \begin{itemize}
        \item MAE y pinball loss de tiempo restante;
        \item AUCPR para riesgo raro;
        \item Brier score y error de calibracion;
        \item cobertura y anchura de intervalos.
      \end{itemize}
    \column{.47\textwidth}
      {\bfseries\color{Navy}Utilidad operativa}\par\smallskip
      \small
      \begin{itemize}
        \item minutos de anticipacion;
        \item falsas alertas por turno/operacion;
        \item cobertura de casos y abstenciones;
        \item utilidad con coste de actuar/no actuar.
      \end{itemize}
  \end{columns}
  \vspace{0.28cm}
  \begin{tikzpicture}
    \node[fill=Navy,rounded corners=5pt,inner sep=9pt,text width=.9\linewidth,text=White] {
      \small\bfseries Objetivo del piloto: demostrar anticipacion accionable, reducir supervision reactiva y construir un business case repetible para terminales y clientes.};
  \end{tikzpicture}
\end{frame}

\begin{frame}{Confianza enterprise como ventaja comercial}
\eyebrow{Adopcion, explicabilidad y control}
  \vspace{0.35cm}
  \begin{columns}[T,onlytextwidth]
    \column{.31\textwidth}
      {\bfseries\color{Blue}Confianza calibrada}\par\smallskip
      \small Cada riesgo incluye intervalo, cobertura y nivel de confianza para priorizar con criterio.
    \column{.31\textwidth}
      {\bfseries\color{Blue}Explicacion operativa}\par\smallskip
      \small Evidencia, cuello de botella y escenario quedan conectados a la operacion y al plan.
    \column{.31\textwidth}
      {\bfseries\color{Blue}Despliegue progresivo}\par\smallskip
      \small Comienza read-only, gana confianza con usuarios y evoluciona hacia recomendaciones integradas.
  \end{columns}
  \vspace{0.65cm}
  \begin{tikzpicture}
    \node[fill=Mint,draw=Cyan!55,rounded corners=5pt,inner sep=8pt,text width=.9\linewidth] {
      \small\bfseries\color{Ink}Auditabilidad y control humano aceleran la adopcion: no son freno, son parte del producto premium.};
  \end{tikzpicture}
\end{frame}

\begin{frame}{Arquitectura EVOCON que acelera el desarrollo}
\eyebrow{Activos tecnicos reutilizables}
  \vspace{0.15cm}
  \small
  \begin{tabularx}{\textwidth}{@{}p{3.35cm}X p{3.25cm}@{}}
    \toprule
    \textbf{Activo} & \textbf{Capacidad aportada} & \textbf{Papel en FlowTwin} \\
    \midrule
    predictiveops-worldmodel & contratos, manifests, splits, leakage tests, gates & nucleo de evidencia \\
    industrial\_jepa\_mvp & agregacion, reporting, estructura experimental & visual analytics \\
    e-jepa-ttc & ventanas, horizontes, incertidumbre, streaming & motor temporal \\
    softnav-jepa & separacion estimador/politica y safety gate & decision segura \\
    SemanticSegmentation\\3Dclouds & disciplina de split y opcion LiDAR futura & expansion multimodal \\
    \bottomrule
  \end{tabularx}
  \vspace{0.35cm}
  \tagbox{Resultado: construimos un nuevo producto sin empezar desde cero}
\end{frame}

\begin{frame}{Plan propuesto: cuatro a seis semanas hasta un piloto}
\eyebrow{De export de datos a demo operativa}
  \vspace{0.35cm}
  \begin{tikzpicture}[x=2.48cm,y=1cm]
    \draw[Slate!35,line width=2pt] (0,0) -- (5,0);
    \foreach \x in {0,...,5} \fill[Cyan] (\x,0) circle (2.6pt);
    \node[above=.2cm,align=center,font=\scriptsize\bfseries,text=Navy] at (0,0) {W0\\Alcance};
    \node[above=.2cm,align=center,font=\scriptsize\bfseries,text=Navy] at (1,0) {W1\\Contrato};
    \node[above=.2cm,align=center,font=\scriptsize\bfseries,text=Navy] at (2,0) {W2\\Proceso};
    \node[above=.2cm,align=center,font=\scriptsize\bfseries,text=Navy] at (3,0) {W3\\Baselines};
    \node[above=.2cm,align=center,font=\scriptsize\bfseries,text=Navy] at (4,0) {W4\\Modelo/sim};
    \node[above=.2cm,align=center,font=\scriptsize\bfseries,text=Navy] at (5,0) {W5--6\\Piloto};
    \node[below=.22cm,align=center,font=\tiny,text=Slate] at (0,0) {decision\\+ coste};
    \node[below=.22cm,align=center,font=\tiny,text=Slate] at (1,0) {timestamps\\+ calidad};
    \node[below=.22cm,align=center,font=\tiny,text=Slate] at (2,0) {mining\\+ conformance};
    \node[below=.22cm,align=center,font=\tiny,text=Slate] at (3,0) {reglas\\+ boosting};
    \node[below=.22cm,align=center,font=\tiny,text=Slate] at (4,0) {PPM/JEPA\\+ escenarios};
    \node[below=.22cm,align=center,font=\tiny,text=Slate] at (5,0) {demo + scorecard\\+ business case};
  \end{tikzpicture}
  \vspace{1.35cm}
  \begin{columns}[T,onlytextwidth]
    \column{.48\textwidth}
      \tagbox{Entregable tangible}\par\smallskip
      \small Dashboard interactivo, API, simulacion y scorecard sobre una operacion seleccionada.
    \column{.48\textwidth}
      \tagbox[Gold!25]{Siguiente expansion}\par\smallskip
      \small Integracion Trace Port y paquete comercial replicable en nuevas terminales y clientes.
  \end{columns}
\end{frame}

\begin{frame}{Como puede generar dinero para Kaleido}
  \eyebrow{Hipotesis a validar, no pricing cerrado}
  \vspace{0.25cm}
  \begin{columns}[T,onlytextwidth]
    \column{.24\textwidth}
      {\bfseries\color{Blue}Add-on SaaS}\par\smallskip
      \small Por terminal, proyecto o volumen de operaciones.
    \column{.24\textwidth}
      {\bfseries\color{Blue}Implantacion}\par\smallskip
      \small Configuracion, mapeo, integracion y calibracion.
    \column{.24\textwidth}
      {\bfseries\color{Blue}White-label}\par\smallskip
      \small Capacidad predictiva dentro de Trace Port.
    \column{.24\textwidth}
      {\bfseries\color{Blue}Uso interno}\par\smallskip
      \small Planificacion, supervision y post-operacion.
  \end{columns}
  \vspace{0.65cm}
  \begin{tikzpicture}
    \node[fill=Navy,rounded corners=6pt,inner sep=11pt,text width=.9\linewidth,text=White,align=center] {
      \bfseries Valor esperado = desviaciones evitadas + horas ahorradas + mejor uso de recursos\\
      $-$ coste de intervenir $-$ coste del modulo};
  \end{tikzpicture}
  \vspace{0.35cm}
  \centering\small\color{Slate}El piloto debe medir los terminos de esta ecuacion con datos de Kaleido.
\end{frame}

\begin{frame}{Aceleradores que ya tenemos}
\eyebrow{No partimos de cero}
  \vspace{0.25cm}
  \begin{columns}[T,onlytextwidth]
    \column{.47\textwidth}
      \small
      \begin{itemize}
        \item Trace Port como sistema de engagement;
        \item acceso a operativas y expertos reales;
        \item productos adyacentes para cross-selling;
        \item datos nuevos generados de forma continua;
        \item canal comercial y marca Kaleido.
      \end{itemize}
    \column{.47\textwidth}
      \small
      \begin{itemize}
        \item contratos, splits y gates ya implementados;
        \item experiencia EVOCON en JEPA y predictive quality;
        \item arquitectura object-centric y multimodelo;
        \item simulacion conectada a decisiones;
        \item roadmap de adaptacion entre terminales.
      \end{itemize}
  \end{columns}
  \vspace{0.45cm}
  \tagbox{La combinacion crea un data moat: cada despliegue mejora el siguiente}
\end{frame}

\begin{frame}{Lo que necesitamos decidir juntos}
  \eyebrow{Siguiente conversacion}
  \vspace{0.35cm}
  \begin{enumerate}
    \item[\color{Cyan}\bfseries 01] \textbf{Operacion:} una carga, descarga o manipulacion repetible con inicio y fin.
    \vspace{0.2cm}
    \item[\color{Cyan}\bfseries 02] \textbf{Decision:} que puede hacer el usuario y con cuanta anticipacion.
    \vspace{0.2cm}
    \item[\color{Cyan}\bfseries 03] \textbf{Evidencia:} esquema, volumen, timestamps y 3--5 casos anonimizados.
    \vspace{0.2cm}
    \item[\color{Cyan}\bfseries 04] \textbf{Economia:} coste de desviacion, recurso e intervencion.
    \vspace{0.2cm}
    \item[\color{Cyan}\bfseries 05] \textbf{Propietarios:} una persona operativa y una tecnica.
  \end{enumerate}
\end{frame}

\begin{frame}{Activemos el piloto FlowTwin}
\eyebrow{Propuesta de arranque}
  \vspace{0.4cm}
  \begin{tikzpicture}
    \node[fill=Navy,rounded corners=8pt,inner sep=16pt,text width=.9\linewidth,text=White] {
      \Large\bfseries Una sesion de 90 minutos, una operacion prioritaria y un
      export inicial para construir la primera demo Kaleido FlowTwin.};
  \end{tikzpicture}
  \vspace{0.65cm}
  \begin{columns}[T,onlytextwidth]
    \column{.31\textwidth}
      {\bfseries\color{Blue}En la primera semana}\par\smallskip
      \small Mapa operativo, contrato de datos y prototipo navegable.
    \column{.31\textwidth}
      {\bfseries\color{Blue}En 4--6 semanas}\par\smallskip
      \small Motor predictivo, simulador, dashboard y scorecard.
    \column{.31\textwidth}
      {\bfseries\color{Blue}Despues}\par\smallskip
      \small Integracion Trace Port y producto premium replicable.
  \end{columns}
\end{frame}

\begin{frame}{Base tecnica consultada}
  \eyebrow{Lectura a 15 julio 2026}
  \vspace{0.15cm}
  \small
  \begin{columns}[T,onlytextwidth]
    \column{.48\textwidth}
      {\bfseries\color{Navy}World models / JEPA}\par\smallskip
      I-JEPA; V-JEPA 2/2.1; LeJEPA; LeWorldModel; Var-JEPA; HWM; FF-JEPA;
      Fast-LeWorldModel; AdaJEPA; WorldDP; SkyJEPA; MIRA; Phys-JEPA y referencias
      asociadas.\par\medskip
      {\bfseries\color{Navy}Aplicacion directa}\par\smallskip
      Predictive Process Monitoring, PGTNet, object-centric graph embeddings,
      OCEL 2.0, survival y calibrated forecasting.
    \column{.48\textwidth}
      {\bfseries\color{Navy}Fuentes de Kaleido}\par\smallskip
      Web corporativa, Kaleido Tech, Trace Port, Shipping Board, Freight
      Intelligence y operaciones portuarias.\par\medskip
      {\bfseries\color{Navy}Auditoria local}\par\smallskip
      Cinco repositorios principales, apuntes JEPA/world models, correo enviado,
      respuesta recibida y propuesta PDF original.
  \end{columns}
  \vspace{0.35cm}
  \centering\tagbox{Bibliografia y trazabilidad completas en el dossier adjunto}
\end{frame}

{\setbeamertemplate{footline}{}
\begin{frame}[plain]
  \begin{tikzpicture}[remember picture,overlay]
    \fill[Navy] (current page.south west) rectangle (current page.north east);
    \fill[Cyan] (current page.south west) rectangle ([xshift=.18cm]current page.north west);
    \fill[Cyan,opacity=.16] ([xshift=13.5cm,yshift=-2cm]current page.north west) circle (3.8cm);
  \end{tikzpicture}
  \vspace{1.25cm}
  {\color{Mint}\bfseries\scriptsize LA OPORTUNIDAD}\par\vspace{0.35cm}
  {\color{White}\fontsize{26}{31}\selectfont\bfseries
  Empezar con foco.\\Escalar con ambicion.\\Convertir operaciones en ventaja.}\par
  \vspace{0.75cm}
  {\color{Mint}\large Kaleido FlowTwin \;--\; inteligencia operativa que mejora con cada operacion}
  \vfill
  {\color{White}\small Álvaro Schwiedop Souto -- Industrial AI \& Predictive Quality Lead, EVOCON Solutions GmbH}
\end{frame}}

\end{document}
